\chapter{Formal Proofs}
\label{appn:A}

This appendix contains the complete proofs of the three theorems stated in
Chapter~\ref{ch:methodology}.

% ============================================================================
\section{Proof of Theorem~1 (Zero-Penalty Guarantee)}
\label{sec:proof_thm1}
% ============================================================================

\noindent\textbf{Theorem~1.} \textit{For any parameter configuration $(f, k)$ and
any biometric sample pair $(\mathcal{M}', \mathbf{c}')$:}
\begin{equation}
  \Pr[\mathrm{Accept}_{\mathrm{IBFV}}] \geq \Pr[\mathrm{Accept}_{\mathrm{Uni}}]
\end{equation}

\noindent\textbf{Proof.}

Let $\mathcal{X}_{\mathrm{FP}}$ be the set of fingerprint-matching vault points
(those whose x-coordinate equals $P_i.x$ for some quantized query minutia $q_i$).
Denote the polynomial degree as $k$, so decoding requires a $(k{+}1)$-tuple of
genuine points for Lagrange interpolation.

The unimodal verification accepts if and only if some $(k{+}1)$-subset of
$\mathcal{X}_{\mathrm{FP}}$ lies on the locked polynomial $p(x)$. Denote this
event $A_\mathrm{Uni}$.

IBFV verification has two exclusive branches:

\textbf{Branch 1: Iris key recovery succeeds.}

Let $\hat{\mathbf{k}}$ be the recovered iris key. The virtual x-coordinates
$\{v_0, \ldots, v_{K-1}\}$ are derived deterministically from $\hat{\mathbf{k}}$
via the PRF in Eq.~(\ref{eq:virtual_x}). Since $v_j = p(v_j)$ was stored at
enrollment (the iris key at enrollment equals the iris key recovered at
verification for a genuine user, so $\hat{\mathbf{k}} = \mathbf{k}$), all $K$
bonus vault points are genuine.

The IBFV candidate set is $\mathcal{X} = \mathcal{X}_{\mathrm{FP}} \cup
\mathcal{X}_{\mathrm{Iris}}$ with $|\mathcal{X}_{\mathrm{Iris}}| = K$.
Smart decoding fixes all $K$ bonus points and enumerates $(k{+}1{-}K)$-subsets
of $\mathcal{X}_{\mathrm{FP}}$.

Formally, IBFV accepts event:
\begin{equation}
  A_\mathrm{IBFV,1} = \{\exists\, S \subseteq \mathcal{X}_{\mathrm{FP}},\,
  |S| = k{+}1{-}K :\; S \cup \mathcal{X}_{\mathrm{Iris}} \text{ lies on } p(x)\}
\end{equation}

Now, $A_\mathrm{Uni}$ occurs iff $\exists\, T \subseteq \mathcal{X}_{\mathrm{FP}},\,
|T| = k{+}1$ lying on $p(x)$. For any such $T$, take $S = T \setminus
\{t_1, \ldots, t_K\}$ (any $K$ elements removed); then $S \cup
\mathcal{X}_{\mathrm{Iris}}$ has $k{+}1$ points on $p(x)$, so $A_\mathrm{IBFV,1}$
also occurs. Therefore $A_\mathrm{Uni} \Rightarrow A_\mathrm{IBFV,1}$, giving:
\begin{equation}
  \Pr[A_\mathrm{IBFV,1}] \geq \Pr[A_\mathrm{Uni}]
\end{equation}

\textbf{Branch 2: Iris key recovery fails.}

IBFV sets $\mathcal{X} = \mathcal{X}_{\mathrm{FP}}$ and applies standard
Lagrange interpolation---identical to unimodal. The acceptance event is:
\begin{equation}
  A_\mathrm{IBFV,2} = A_\mathrm{Uni}
\end{equation}
and trivially $\Pr[A_\mathrm{IBFV,2}] = \Pr[A_\mathrm{Uni}]$.

\textbf{Combining:}

Let $p_s = \Pr[\text{iris recovery succeeds}]$. By the law of total probability:
\begin{align}
  \Pr[A_\mathrm{IBFV}] &= p_s \cdot \Pr[A_\mathrm{IBFV,1}] +
  (1{-}p_s) \cdot \Pr[A_\mathrm{IBFV,2}] \\
  &\geq p_s \cdot \Pr[A_\mathrm{Uni}] + (1{-}p_s) \cdot \Pr[A_\mathrm{Uni}] \\
  &= \Pr[A_\mathrm{Uni}]
\end{align}
\hfill$\square$

% ============================================================================
\section{Proof of Theorem~2 (Share Confidentiality)}
\label{sec:proof_thm2}
% ============================================================================

\noindent\textbf{Theorem~2.} \textit{For any $t' < t$ compromised shares, the
conditional distribution of vault $\mathbf{b}$ is uniform over
$\{0, \ldots, 255\}^L$. Fewer than $t$ shares reveal zero information about the
vault content.}

\noindent\textbf{Proof.}

Fix any byte index $i$. The Shamir polynomial for byte $i$ is:
\begin{equation}
  f_i(x) = b_i + a_1^{(i)} x + \cdots + a_{t-1}^{(i)} x^{t-1} \pmod{257}
\end{equation}
where coefficients $a_j^{(i)} \stackrel{\$}{\leftarrow} \mathrm{GF}(257)$
are chosen uniformly and independently at random.

Suppose the adversary holds shares at positions $x_1, \ldots, x_{t'}$ with
$t' < t$. This gives the system of equations:
\begin{equation}
  \begin{pmatrix} f_i(x_1) \\ \vdots \\ f_i(x_{t'}) \end{pmatrix}
  = \underbrace{\begin{pmatrix} 1 & x_1 & x_1^2 & \cdots & x_1^{t-1} \\
  \vdots & \vdots & \vdots & & \vdots \\
  1 & x_{t'} & x_{t'}^2 & \cdots & x_{t'}^{t-1}
  \end{pmatrix}}_{\mathbf{V} \in \mathrm{GF}(257)^{t' \times t}}
  \begin{pmatrix} b_i \\ a_1^{(i)} \\ \vdots \\ a_{t-1}^{(i)} \end{pmatrix}
\end{equation}

$\mathbf{V}$ is a Vandermonde matrix with $t' < t$ rows and $t$ columns over
$\mathrm{GF}(257)$ (a field). Since $t' < t$, the system is underdetermined:
for any fixed observed share values and any candidate $b_i \in \{0, \ldots, 255\}$,
there exists at least one choice of the remaining coefficients
$(a_1^{(i)}, \ldots, a_{t-1}^{(i)})$ that satisfies the system. The mapping
from coefficients to shares is linear and surjective from a uniform random
distribution over $\mathrm{GF}(257)^t$, which projects uniformly onto any $t'$
shares. Therefore, the conditional distribution of $b_i$ given
$t' < t$ shares is the uniform distribution over $\mathrm{GF}(257) \supseteq
\{0, \ldots, 255\}$.

Since bytes are shared independently, the joint distribution of $\mathbf{b}$
given $t' < t$ shares is uniform over $\mathrm{GF}(257)^L \supseteq
\{0,\ldots,255\}^L$. This is information-theoretic security: the mutual
information $I(\mathbf{b}; \text{shares}_{< t}) = 0$.
\hfill$\square$

% ============================================================================
\section{Proof of Theorem~3 (Exact Reconstruction)}
\label{sec:proof_thm3}
% ============================================================================

\noindent\textbf{Theorem~3.} \textit{Given any $t$ shares with distinct indices,
Lagrange interpolation over $\mathrm{GF}(257)$ recovers $\hat{b}_i = b_i$ for
all $i$, yielding $\hat{\mathbf{b}} = \mathbf{b}$ with zero bit errors.}

\noindent\textbf{Proof.}

Fix byte $i$. The polynomial $f_i(x)$ has degree $t{-}1$ over $\mathrm{GF}(257)$.
By the fundamental theorem of algebra over finite fields, a polynomial of degree
$t{-}1$ is uniquely determined by its values at any $t$ distinct points.

Given shares $(x_1, f_i(x_1)), \ldots, (x_t, f_i(x_t))$ at distinct
$x_1, \ldots, x_t \in \mathrm{GF}(257)$, Lagrange interpolation reconstructs:
\begin{equation}
  \hat{f}_i(x) = \sum_{j=1}^{t} f_i(x_j)
  \prod_{\substack{\ell=1 \\ \ell \neq j}}^{t}
  \frac{x - x_\ell}{x_j - x_\ell}
\end{equation}
where all arithmetic is in $\mathrm{GF}(257)$. The interpolated polynomial
$\hat{f}_i$ is the unique degree-$(t{-}1)$ polynomial through the $t$ given
points, hence $\hat{f}_i \equiv f_i$. Evaluating at $x = 0$:
\begin{equation}
  \hat{b}_i = \hat{f}_i(0) = f_i(0) = b_i
\end{equation}
exactly. Since this holds for each byte independently and arithmetic is exact
over the finite field, $\hat{\mathbf{b}} = \mathbf{b}$ with zero bit errors.
\hfill$\square$
